% TEMPLATE-MILC-v3.tex - plantilla maestra UNICA de las guias MILC v3.
%
% La usan las tres familias de guias, cada una con su driver:
%   · Grados 6-11  -> scripts/build-guias-g11.py
%   · Bebras       -> scripts/build-guias-bebras.py
%   · Semillero    -> scripts/build-guias-semillero.py
%
% Antes habia tres copias casi identicas (~400 lineas iguales, 6 distintas):
% cada mejora habia que aplicarla tres veces, y en la practica no se hacia
% (el motor de recursos vivia solo en dos de ellas). Lo que varia entre
% familias esta parametrizado en 6 placeholders que cada driver llena
% (se nombran aqui SIN su sintaxis real: el builder reemplaza por texto plano,
% tambien dentro de los comentarios, y un valor multilinea rompe el preambulo):
%
%   HEADER_IZQ        encabezado izquierdo de pagina
%   HEADER_DER        encabezado derecho de pagina
%   PORTADA_ETIQUETA  rotulo sobre el numero grande (GRADO/MOMENTO/MODULO)
%   PORTADA_SERIE     linea de serie de la portada
%   PORTADA_ID        identificador de la guia en la portada
%   PORTADA_CIRCULO   el circulito de grado (°), o vacio si no aplica
%
% El resto de claves las documenta content/guias/_SCHEMA.md. El script valida
% que no queden placeholders sin reemplazar antes de escribir el archivo final.

\documentclass[11pt,letterpaper]{article}
\usepackage[margin=2.0cm]{geometry}
\usepackage[table]{xcolor}
\usepackage{fontspec}
\usepackage[spanish]{babel}
\usepackage{tikz}
% Librerías TikZ para los diagramas de las guías (bloque `recursos:`):
%  · babel  → babel en español vuelve activos caracteres como «>», y TikZ los usa
%             en sus opciones (>=stealth, ->). Sin esta librería, cualquier
%             diagrama con flechas revienta con «\language@active@arg> has an
%             extra }». Debe ir ANTES de que se dibuje nada.
%  · positioning        → sintaxis «below left=2pt and 3pt».
%  · decorations.pathreplacing → llaves ({brace}) para señalar rangos.
\usetikzlibrary{babel,positioning,decorations.pathreplacing}
\usepackage{graphicx}
% Circuitos: el trabajo de electrónica se hace hoy con micro:bit y sus sensores
% integrados (MakeCode, sin cableado), así que no se necesitan esquemáticos.
% Si algún día entran montajes con protoboard, basta descomentar esta línea:
% los diagramas .tex/.tikz declarados en `recursos:` ya se incrustan solos.
% \usepackage{circuitikz}
\usepackage[most]{tcolorbox}
\usepackage{tabularx}
\usepackage{array}
\usepackage{fancyhdr}
\usepackage{titlesec}
\usepackage[document]{ragged2e}  % bandera a la derecha en todo el cuerpo
\usepackage{enumitem}
\usepackage{microtype}

% Helvetica Neue no trae la flecha U+2192, asi que cada «→» escrito literalmente
% en el YAML se compilaba a NADA, en silencio (462 flechas en 61 guias: rutas del
% tipo «paleta → tipografia → lockup» salian rotas). Mapeamos el caracter a su
% version matematica, que si tiene glifo. Asi el autor puede seguir escribiendo
% «→» comodamente en el YAML.
\usepackage{newunicodechar}
\newunicodechar{→}{\ensuremath{\rightarrow}}
% Mismo problema con los simbolos de marca y vinetas: el «✓» de «una palomita ✓»
% salia como cuadrito vacio en el PDF de todas las guias que lo usaban.
\usepackage{pifont}
\newunicodechar{✓}{\ding{51}}
\newunicodechar{✗}{\ding{55}}
\newunicodechar{✔}{\ding{52}}
\newunicodechar{•}{\textbullet}
\newunicodechar{≥}{\ensuremath{\geq}}
\newunicodechar{≤}{\ensuremath{\leq}}

\setmainfont{Helvetica Neue}
\setsansfont{Helvetica Neue}
\setlength{\parindent}{0pt}
\setlength{\parskip}{6pt}
% Sin particion de palabras: en 6.º-7.º una palabra cortada por guion al final
% de linea («Comuni-cacion») frena la lectura mucho mas que una linea corta.
\hyphenpenalty=10000
\exhyphenpenalty=10000
\pretolerance=2000
\tolerance=3000
\setlength{\emergencystretch}{3em}
\linespread{1.10}
\renewcommand{\arraystretch}{1.25}
\setlist{leftmargin=5mm,itemsep=1mm,topsep=1mm}
\newcolumntype{Y}{>{\RaggedRight\arraybackslash}X}

\definecolor{milcMagenta}{HTML}{E6007E}
\definecolor{milcVino}{HTML}{5A0038}
\definecolor{milcTurquesa}{HTML}{0093A5}
\definecolor{milcVerde}{HTML}{58A923}
\definecolor{milcMostaza}{HTML}{E5B400}
\definecolor{milcOcre}{HTML}{B86600}
\definecolor{milcNegro}{HTML}{241816}
\definecolor{milcGris}{HTML}{F7F7F4}
\definecolor{milcLinea}{HTML}{E8E2DD}
\definecolor{milcGradePrimary}{HTML}{FF2D87}
\definecolor{milcGradeDark}{HTML}{9F1B56}
\definecolor{milcGradeSoft}{HTML}{FFE0EE}
\definecolor{milcGradeAccent}{HTML}{CFE7FF}
\definecolor{milcGradeText}{HTML}{FFFFFF}
\color{milcNegro}

\pagestyle{fancy}
\fancyhf{}
\lhead{\small Tecnología e Informática · Grado 10}
\rhead{\small Guía 19 · MILC v3}
\cfoot{\small\thepage}
\renewcommand{\headrulewidth}{0pt}

\newtcolorbox{titlebox}[1]{
  enhanced,colback=#1,colframe=#1,arc=7mm,boxrule=0pt,
  left=7mm,right=7mm,top=3mm,bottom=3mm,
  before skip=8pt,after skip=8pt,
  fontupper=\sffamily\bfseries\Large\color{white}
}

\newtcolorbox{softbox}[3]{
  enhanced,colback=#2,colframe=#1,borderline west={4pt}{0pt}{#1},
  arc=2mm,boxrule=.4pt,left=4mm,right=4mm,top=3mm,bottom=3mm,
  before skip=6pt,after skip=8pt,title={#3},coltitle=white,
  colbacktitle=#1,fonttitle=\sffamily\bfseries,fontupper=\RaggedRight,
  attach boxed title to top left={xshift=3mm,yshift=-2mm},
  boxed title style={arc=2mm, boxrule=0pt}
}

\newtcolorbox{drawbox}[2]{
  enhanced,colback=white,colframe=#1,arc=4mm,boxrule=1pt,
  left=5mm,right=5mm,top=4mm,bottom=4mm,before skip=8pt,after skip=8pt,
  title={#2},coltitle=white,colbacktitle=#1,fonttitle=\sffamily\bfseries,
  fontupper=\RaggedRight,
  attach boxed title to top left={xshift=4mm,yshift=-2mm},
  boxed title style={arc=3mm, boxrule=0pt}
}

\newtcolorbox{infoband}[1]{
  enhanced,colback=#1!8,colframe=#1!50,arc=2mm,boxrule=.5pt,
  left=4mm,right=4mm,top=2mm,bottom=2mm,before skip=4pt,after skip=4pt,
  fontupper=\small\RaggedRight
}

% Puente narrativo entre secciones (contrato editorial Conéctate).
% Da continuidad: "Acabas de X. Ahora vas a Y porque Z."
% Visualmente: banda izquierda en color del grado, texto en itálica pequeña.
\newtcolorbox{puentebox}{
  enhanced,colback=milcGradeSoft,colframe=milcGradeDark,
  borderline west={3pt}{0pt}{milcGradeDark},
  arc=0mm,boxrule=0pt,left=5mm,right=4mm,top=2.5mm,bottom=2.5mm,
  before skip=4pt,after skip=8pt,
  fontupper=\itshape\small\color{milcNegro}\RaggedRight
}
% La flecha va en modo matematico a proposito: Helvetica Neue NO tiene el glifo
% U+2192, asi que \textrightarrow se compilaba a NADA («Missing character: There
% is no → in font Helvetica Neue») y el puente salia sin flecha. En modo
% matematico la flecha viene de la fuente matematica y si se dibuja.
\newcommand{\puente}[1]{%
  \begin{puentebox}%
    \textcolor{milcGradeDark}{$\rightarrow$}\hspace{2mm}#1%
  \end{puentebox}%
}

\newcommand{\linead}[1]{\noindent\color{milcLinea}\rule{#1}{0.5pt}\color{milcNegro}}
\newcommand{\respuesta}[1]{\par\vspace{2mm}\foreach \n in {1,...,#1}{\noindent\color{milcLinea}\rule{\linewidth}{0.5pt}\par\vspace{5mm}}\color{milcNegro}}
\newcommand{\checkbox}{\textcolor{milcGradeDark}{\fbox{\phantom{x}}}\hspace{5pt}}

\newcommand{\phasecard}[4]{
\begin{tcolorbox}[enhanced,colback=#3,colframe=#2,arc=3mm,boxrule=.5pt,height=3.05cm,valign=center,halign=center,left=1.5mm,right=1.5mm,top=2mm,bottom=2mm]
{\sffamily\bfseries\fontsize{22}{24}\selectfont\textcolor{#2}{#1}}\par
{\sffamily\bfseries\small #4}
\end{tcolorbox}
}

% ── Piezas del rediseno 2026 (legibilidad 6.º-11.º) ───────────────────
% Banda de estacion: reemplaza el titlebox macizo. Titulo a la izquierda,
% dato practico (tiempo/modalidad) a la derecha, en una sola linea.
\newcommand{\milcestacion}[2]{%
\begin{tcolorbox}[enhanced,colback=#1,colframe=#1,arc=2mm,boxrule=0pt,
  left=5mm,right=5mm,top=2.6mm,bottom=2.6mm,before skip=11pt,after skip=7pt]
{\sffamily\bfseries\fontsize{15}{18}\selectfont\color{white}#2}
\end{tcolorbox}}

% Tarjeta de paso numerada: sustituye las tablas «Pilar / Aplicacion», que
% eran formato de consulta docente, no de estudiante. El numero va en un
% circulo sobre el borde para que la secuencia se lea de un vistazo.
\newcommand{\milcpaso}[3]{%
\begin{tcolorbox}[enhanced,colback=white,colframe=milcLinea,arc=1.5mm,boxrule=.9pt,
  left=14mm,right=4mm,top=3mm,bottom=3mm,before skip=4pt,after skip=4pt,
  fontupper=\RaggedRight,
  overlay={\node[anchor=north west,circle,fill=#2,text=white,inner sep=0pt,
    minimum size=7.4mm,font=\sffamily\bfseries\small]
    at ([xshift=3.6mm,yshift=-3.2mm]frame.north west) {#1};}]
#3
\end{tcolorbox}}

% Casilla marcable de tamano util con lapiz (la anterior era \fbox{\phantom{x}},
% de alto variable segun el texto que la seguia).
\newcommand{\milcmarca}{\framebox[3.8mm]{\rule{0pt}{3.4mm}}}

% Fila marcable de lista suelta (checks de la estacion 1).
\newcommand{\milccheckrow}[1]{%
\par\vspace{1.8mm}\noindent
\makebox[6.4mm][l]{\raisebox{-.55mm}{\textcolor{milcNegro}{\milcmarca}}}%
\parbox[t]{\dimexpr\linewidth-6.4mm\relax}{\RaggedRight #1}\par}

% Celda marcable dentro de la rubrica (tabla de autoevaluacion). Misma
% sangria francesa, pero pensada para vivir dentro de una columna X.
\newcommand{\milcopcion}[1]{%
\makebox[5.2mm][l]{\raisebox{-.55mm}{\textcolor{milcNegro}{\milcmarca}}}%
\parbox[t]{\dimexpr\linewidth-5.2mm\relax}{\RaggedRight #1}}

% ── Recursos visuales de la guía ──────────────────────────────────────
% Declarados en el bloque `recursos:` del YAML y emitidos por el builder.
% \guiaFigura   → imagen rasterizada o PDF (foto, carta celeste, montaje).
% \guiaDiagrama → fuente TikZ/circuitikz (.tex) incrustada como vector:
%                 es la vía para circuitos y esquemáticos editables.
% #1 = ruta relativa al .tex · #2 = pie (puede ir vacío)
\newcommand{\guiaFigura}[2]{%
\begin{center}
\includegraphics[width=0.88\linewidth,height=0.38\textheight,keepaspectratio]{#1}\\[1.5mm]
{\footnotesize\itshape\textcolor{milcNegro!75}{#2}}
\end{center}
}

\newcommand{\guiaDiagrama}[2]{%
\begin{center}
\input{#1}\\[1.5mm]
{\footnotesize\itshape\textcolor{milcNegro!75}{#2}}
\end{center}
}

\begin{document}
\thispagestyle{empty}

% ── Portada ───────────────────────────────────────────────────────────
% Antes: un rectangulo de color a pagina completa. Se veia bien en pantalla,
% pero en la fotocopiadora del colegio se comia el toner y salia gris sucio.
% Ahora la banda ocupa el tercio superior y el resto va en blanco.
\begin{tikzpicture}[remember picture,overlay]
  \path[fill=milcGradePrimary]
    (current page.north west) rectangle ([yshift=-10.6cm]current page.north east);

  \node[anchor=north west,text=milcGradeText] at ([xshift=2.0cm,yshift=-1.55cm]current page.north west)
    {\sffamily\bfseries\fontsize{12}{14}\selectfont GRADO};
  \node[anchor=north west,text=milcGradeText,opacity=.85] at ([xshift=2.0cm,yshift=-2.35cm]current page.north west)
    {\sffamily\bfseries\fontsize{8.5}{10}\selectfont SERIE GUÍAS · TECNOLOGÍA E INFORMÁTICA};
  \node[anchor=north east,text=milcGradeText,opacity=.85,align=right] at ([xshift=-2.0cm,yshift=-1.55cm]current page.north east)
    {\sffamily\bfseries\fontsize{9}{12}\selectfont I.E. Sor María Juliana\\Cartago, Valle del Cauca};

  \node[anchor=west,text=milcGradeText] (gradonum) at ([xshift=1.85cm,yshift=-7.1cm]current page.north west)
    {\sffamily\bfseries\fontsize{112}{112}\selectfont 10};
  % El circulito de grado (°) solo aplica a las guias de grado («11°»); en
  % Bebras y Semillero el numero grande es un momento/modulo, no un grado.
  % Se posiciona relativo al nodo `gradonum`, no en coordenadas absolutas.
  \draw[milcGradeText,line width=1.6pt]
    ([xshift=2.0mm,yshift=-4.5mm]gradonum.north east) circle (.15cm);

  \node[anchor=west,text=milcGradeText,text width=12.3cm,align=left] at ([xshift=6.5cm,yshift=-6.6cm]current page.north west)
    {
     {\sffamily\bfseries\fontsize{9}{11}\selectfont\MakeUppercase{Período 2 · Informes técnicos y comunicación profesional con IA}}\\[3mm]
     {\sffamily\bfseries\fontsize{21}{25}\selectfont Mini-proyecto P2 ---\\informe técnico\\con datos reales}\\[3.5mm]
     {\sffamily\bfseries\fontsize{15}{18}\selectfont Guía 19-10-TIC}
    };
\end{tikzpicture}

\vspace*{9.4cm}
\vfill

{\sffamily\bfseries\fontsize{8}{10}\selectfont\textcolor{milcGradeDark}{LO QUE TE LLEVAS HOY}}

\begin{tcolorbox}[enhanced,colback=white,colframe=milcNegro,arc=2mm,boxrule=1.6pt,
  left=5mm,right=5mm,top=4mm,bottom=4mm,before skip=3pt,after skip=12pt,
  fontupper=\RaggedRight\fontsize{11}{15.5}\selectfont]
Mini-proyecto integrador del periodo 2: informe técnico de 6-10 páginas sobre un tema escolar real con (a) datos reales recolectados con encuesta de la sesión 5 más análisis cualitativo de la sesión 6, (b) estructura completa del informe técnico de la sesión 2 (introducción, desarrollo, conclusiones), (c) maquetación profesional en Google Docs o Markdown (sesiones 3-4), (d) declaración explícita y honesta de uso de IA con prompts, porcentajes y bitácora, (e) presupuesto si aplica + carta del autor firmada de 1 página
\end{tcolorbox}

\begin{tcolorbox}[enhanced,colback=milcGris,colframe=milcGris,arc=2mm,boxrule=0pt,
  left=5mm,right=5mm,top=3.5mm,bottom=3.5mm,after skip=12pt]
{\sffamily\bfseries\fontsize{8}{10}\selectfont\textcolor{milcNegro!55}{ESTA GUÍA ES DE}}\\[3mm]
\begin{tabularx}{\linewidth}{X p{3.2cm} p{3.2cm}}
\linead{\linewidth} & \linead{3.0cm} & \linead{3.0cm}\\[-1mm]
{\fontsize{8.5}{10}\selectfont\textcolor{milcNegro!55}{Nombre completo}} &
{\fontsize{8.5}{10}\selectfont\textcolor{milcNegro!55}{Grupo}} &
{\fontsize{8.5}{10}\selectfont\textcolor{milcNegro!55}{Fecha}}\\
\end{tabularx}
\end{tcolorbox}

\vfill

\noindent\textcolor{milcLinea}{\rule{\linewidth}{1pt}}\\[2mm]
{\sffamily\bfseries\fontsize{9.5}{12}\selectfont Autor: Dr. Álvaro Cárdenas Orozco}\\[1mm]
{\sffamily\fontsize{8.5}{11}\selectfont\textcolor{milcNegro!55}{Metodología MILC · apertura ancestral · pensamiento computacional · triángulo Dussel–estoicismo–Floridi}}

\newpage

\section*{Apertura: Mini-proyecto --- informe técnico con datos reales y IA documentada}

\begin{softbox}{milcGradeDark}{milcGradeSoft}{Saber ancestral}
En los gremios tradicionales del oficio (carpinteros, ebanistas, sastres, alfareros) hubo durante siglos una práctica formal que cualquier maestro conocía: \textbf{el aprendiz que pedía certificación al gremio debía presentar una pieza maestra}. No cualquier pieza: una específica que demostrara que el aprendiz dominaba \textbf{todas las técnicas aprendidas en el periodo de formación}. Si el aprendiz era ebanista, la pieza debía integrar: corte de madera, ensamblaje sin clavos, lijado en 3 grados, terminación con barniz, decoración tallada. Si saltaba una técnica, la pieza no era certificable. El maestro y los oficiales viejos del gremio la examinaban con detalle: pasaban la mano por las uniones, golpeaban la madera para escuchar resonancia, miraban a contraluz el lijado. Si la pieza integraba todas las técnicas con calidad, el aprendiz recibía el \textbf{título de oficial}. Si fallaba en alguna, volvía al taller a perfeccionar la técnica faltante. La sabiduría del gremio era inquebrantable: \textbf{el oficio se certifica con pieza integradora, no con conocimientos sueltos}. El mini-proyecto del periodo 2 es esa pieza maestra moderna: integra las 8 sesiones anteriores en un solo informe técnico profesional.
\end{softbox}

\begin{softbox}{milcTurquesa}{milcGris}{Saber contemporáneo}
El \textbf{mini-proyecto integrador del periodo 2} es una pieza única que articula todo lo aprendido: distinción de tipos de texto (S1), estructura del informe técnico (S2), Google Docs profesional (S3), Markdown si conviene (S4), encuestas con Forms más IA (S5), análisis cualitativo con auditoría humana (S6), opcionalmente lo persuasivo del informe comercial (S7) y la comunicación profesional por correo (S8). Sus componentes obligatorios: (1) \textbf{Tema escolar real} con destinatario potencial. (2) \textbf{Datos reales} recolectados con tu encuesta (15-20 respuestas mínimas) más análisis cualitativo auditado. (3) \textbf{Estructura completa de informe técnico}: introducción (contexto + objetivo + alcance), desarrollo (metodología + datos + análisis), conclusiones (hallazgos + recomendaciones). (4) \textbf{Maquetación profesional} en Google Docs con estilos, TOC, citas APA, imagen con pie. (5) \textbf{Declaración honesta de IA} con bitácora detallada: qué modelos usaste, qué prompts, qué porcentaje aproximado de texto generaste con IA, qué editaste a mano. (6) \textbf{Carta del autor firmada} (1 página): declaración personal del proceso y compromiso con la propuesta.
\end{softbox}

\begin{softbox}{milcMostaza}{milcGris}{Pregunta-puente}
\textit{¿Qué sabía el aprendiz del gremio al presentar pieza maestra que integraba todas las técnicas, que el estudiante novato olvida cuando entrega trabajos sueltos sin integración? ¿Y por qué la carta del autor firmada es lo que separa el informe profesional del ejercicio escolar?}
\end{softbox}

\begin{softbox}{milcVerde}{milcGris}{Saber y saber hacer}
Al terminar podrás: (1) \textbf{analizar} el inventario de todo lo que produjiste en las 8 sesiones del periodo y planificar cómo integrarlo; (2) \textbf{aplicar} los hallazgos de tu encuesta y análisis a un informe técnico completo dirigido a destinatario real; (3) \textbf{crear} el mini-proyecto con maquetación profesional y declaración honesta de uso de IA; (4) \textbf{evaluar} tu informe con checklist de revisión antes de declararlo final.
\end{softbox}



\small
\begin{tabularx}{\textwidth}{>{\bfseries}p{3.0cm}Y}
\rowcolor{milcGradePrimary!12}Autor & Dr. Álvaro Cárdenas Orozco\\
DBA & Comunicación profesional con asistencia de IA y herramientas ofimáticas gratuitas (MEN, T\&I)\\
Referentes & Markdown como estándar abierto · Google Workspace · Estoicismo (claridad)\\
Producto final & Mini-proyecto integrador del periodo 2: informe técnico de 6-10 páginas sobre un tema escolar real con (a) datos reales recolectados con encuesta de la sesión 5 más análisis cualitativo de la sesión 6, (b) estructura completa del informe técnico de la sesión 2 (introducción, desarrollo, conclusiones), (c) maquetación profesional en Google Docs o Markdown (sesiones 3-4), (d) declaración explícita y honesta de uso de IA con prompts, porcentajes y bitácora, (e) presupuesto si aplica + carta del autor firmada de 1 página\\
\end{tabularx}
\normalsize

\puente{Antes de empezar, mira el viaje completo. Has recorrido 8 sesiones: tipos de texto, estructura, Docs, Markdown, encuestas, análisis IA, informe comercial, correo profesional. Hoy es el \textbf{día de integrar}: las 8 piezas en un solo informe técnico profesional. La sesión 10 será sustentación pública del mini-proyecto.}

\milcestacion{milcGradeDark}{Ruta MILC: así vamos a aprender}
\begin{tabularx}{\textwidth}{>{\centering\arraybackslash}X>{\centering\arraybackslash}X>{\centering\arraybackslash}X>{\centering\arraybackslash}X}
\phasecard{1}{milcVerde}{milcGris}{Escucha\\[-1mm]\small Observo, escucho y pregunto.} &
\phasecard{2}{milcTurquesa}{milcGris}{Entender\\[-1mm]\small Sistematizo y ordeno.} &
\phasecard{3}{milcMagenta}{milcGris}{Hacer\\[-1mm]\small Construyo y aplico.} &
\phasecard{4}{milcVino}{milcGris}{Cierre\\[-1mm]\small Evalúo y reflexiono.}
\end{tabularx}


\puente{Antes de teorizar, vas a hacer un \emph{inventario completo}: lista todo lo que tienes producido en el periodo (datos de encuesta de S5, análisis auditado de S6, informe comercial de S7 si aplica, plantilla de Docs de S3 con estilos aplicados) y lo que aún falta para el mini-proyecto.}

\milcestacion{milcVerde}{Estación 1 de 3 · Escucha}

\begin{softbox}{milcVerde}{milcGris}{Escena para escuchar}
\textbf{Actividad 1 · ANALIZA --- Inventario completo del periodo} (15 min · individual). Haz lista de lo que tienes producido en las 8 sesiones del periodo. Reglas: (1) \textbf{Tengo}: plantilla de informe técnico (S2), datos de encuesta con 15-20 respuestas (S5), análisis cualitativo auditado (S6), plantilla de Docs con estilos aplicados (S3), tal vez informe comercial completo (S7). (2) \textbf{Me falta}: integración en informe técnico completo, declaración detallada de uso de IA, carta del autor, revisión final. Para cada item faltante, estima tiempo necesario.
\end{softbox}

{\sffamily\bfseries\fontsize{8}{10}\selectfont\textcolor{milcNegro!55}{MARCA CUANDO LO HAYAS HECHO}}
\milccheckrow{Listé lo producido en las 8 sesiones.}
\milccheckrow{Listé lo que falta para el mini-proyecto.}
\milccheckrow{Estimé tiempo necesario para cerrar pendientes.}

\begin{infoband}{milcVerde}


\textbf{Lo que va al cuaderno (Actividad 1 · ANALIZA).} Título: «Actividad 1 --- Inventario integración». \textbf{Formato:} 2 listas (tengo, me falta) + cronograma esta semana. \textbf{Sabes que terminaste cuando} tienes plan claro para integrar.
\end{infoband}

\puente{Ya tienes inventario. Ahora vas a entender la \textbf{anatomía del mini-proyecto integrador}: 6 componentes, los 5 errores típicos del aprendiz que se apura, la cadena lógica entre datos y recomendaciones.}

\milcestacion{milcTurquesa}{Estación 2 de 3 · Entender}

\begin{softbox}{milcTurquesa}{milcGris}{Lo que vas a entender}
Tu inventario es el plan. Ahora necesitas la \textbf{anatomía del mini-proyecto integrador} para que la pieza maestra funcione.
\end{softbox}

\milcpaso{1}{milcTurquesa}{Los \textbf{6 componentes obligatorios} se descomponen así: (1) \textbf{Tema y destinatario}: declarar al inicio el tema escolar real y el destinatario potencial (rector, coordinador, JAC, microempresario familiar). (2) \textbf{Datos reales}: anexo o sección con tu encuesta aplicada y los 15-20+ resultados. (3) \textbf{Estructura informe técnico}: introducción (contexto/objetivo/alcance), desarrollo (metodología/datos/análisis), conclusiones (hallazgos/recomendaciones accionables). 6-10 páginas en total. (4) \textbf{Maquetación profesional}: Google Docs con estilos de título, TOC autogenerada, numeración, citas APA, imagen con pie de figura. (5) \textbf{Declaración de IA}: bitácora detallada con prompts, modelos, porcentajes, qué editaste tú. (6) \textbf{Carta del autor firmada}: 1 página con tu compromiso personal, proceso seguido, qué aprendiste.}
\milcpaso{2}{milcTurquesa}{La \textbf{cadena lógica} es crucial: los datos del desarrollo deben sostener las conclusiones; las conclusiones deben generar recomendaciones; las recomendaciones deben tener destinatario real con capacidad de implementarlas. Si rompes la cadena (datos que no sostienen conclusiones, recomendaciones sin responsable), el informe pierde credibilidad. La phronesis del oficio consiste en \textbf{verificar que cada eslabón sostenga al siguiente}.}
\milcpaso{3}{milcTurquesa}{Los \textbf{5 errores típicos} del mini-proyecto novato: (a) \textbf{Datos insuficientes}: menos de 15 respuestas no permiten análisis. (b) \textbf{Análisis sin auditar la IA}: aceptar clasificaciones sin verificar. (c) \textbf{Recomendaciones sin responsables}: \emph{``alguien debería hacer X''}. (d) \textbf{Sin maquetación profesional}: texto crudo sin estilos. (e) \textbf{Declaración de IA vaga}: \emph{``usé IA''} sin detalles.}
\milcpaso{4}{milcTurquesa}{Algoritmo: (1) inventario de lo producido. (2) integrar datos en una sola sección. (3) escribir introducción según S2. (4) desarrollo con datos reales y análisis auditado. (5) conclusiones con recomendaciones accionables. (6) maquetar en Docs con S3. (7) declarar IA en bitácora detallada. (8) escribir carta del autor firmada. (9) revisión final con checklist.}

\begin{drawbox}{milcTurquesa}{Anatomía del mini-proyecto}
\textbf{Tema y destinatario:} reales nombrados. \\[1mm] \textbf{Datos reales:} 15-20 respuestas mínimas. \\[1mm] \textbf{Estructura informe:} intro + desarrollo + conclusiones. \\[1mm] \textbf{Maquetación:} Docs profesional. \\[1mm] \textbf{Declaración IA:} bitácora detallada. \\[1mm] \textbf{Carta autor:} 1 página firmada.
\end{drawbox}

\begin{softbox}{milcMostaza}{milcGris}{Errores comunes y su corrección}
(1) Datos insuficientes (menos de 15). (2) Análisis sin auditar IA. (3) Recomendaciones sin responsables. (4) Sin maquetación profesional. (5) Declaración de IA vaga. (6) Sin carta del autor. (7) Cadena lógica rota.
\end{softbox}

\begin{infoband}{milcTurquesa}


\textbf{Lo que va al cuaderno (Actividad 2 · APLICA).} Título: «Actividad 2 --- Anatomía mini-proyecto». \textbf{Formato:} (a) los 6 componentes; (b) cadena lógica datos-conclusiones-recomendaciones; (c) los 7 errores con marca. \textbf{Sabes que terminaste cuando} tienes plan claro de integración.
\end{infoband}

\puente{Tienes el método. Toca aplicarlo: integras componentes, escribes el informe completo, lo maquetas en Docs, llenas bitácora de IA, escribes carta del autor firmada.}

\milcestacion{milcMagenta}{Estación 3 de 3 · Hacer}

\begin{softbox}{milcMagenta}{milcGris}{Cómo vas a construir tu producto}
\textbf{Actividad 3 · CREA + EVALÚA --- Mini-proyecto completo} (60 min en sesión + tiempo en casa · individual). Hora de aplicar. Integras componentes, escribes informe completo, maquetas profesional, declaras IA, escribes carta.
\end{softbox}

\milcpaso{1}{milcMagenta}{Tu mini-proyecto se completa en 8 pasos: (1) abrir documento en Google Docs con plantilla profesional (S3). (2) escribir introducción con contexto/objetivo/alcance/destinatario. (3) escribir desarrollo integrando metodología (cómo aplicaste la encuesta), datos (los 15-20+ resultados), análisis (de S6 con auditoría). (4) escribir conclusiones con hallazgos y recomendaciones accionables con responsables. (5) generar TOC, agregar citas APA, insertar imagen con pie de figura. (6) escribir bitácora detallada de uso de IA. (7) escribir carta del autor firmada de 1 página. (8) revisión final.}
\milcpaso{2}{milcMagenta}{Sigue el patrón \textbf{``integración honesta sobre yuxtaposición de piezas''}: el mini-proyecto no es solo poner las 8 piezas en orden; es \emph{tejerlas en un solo argumento}. La introducción presenta el problema, el desarrollo lo investiga con datos, las conclusiones proponen acción. Cada sesión anterior se transforma en una sección del mini-proyecto.}
\milcpaso{3}{milcMagenta}{La \textbf{carta del autor firmada} tiene 4 párrafos: (a) \textbf{Qué}: el mini-proyecto en 2-3 líneas. (b) \textbf{Por qué}: por qué este tema, por qué este destinatario. (c) \textbf{Proceso}: cómo lo hiciste con honestidad sobre uso de IA. (d) \textbf{Compromiso}: qué te llevas como aprendiz. Cierra con nombre completo + fecha + ciudad.}
\milcpaso{4}{milcMagenta}{Algoritmo: (1) Docs con plantilla. (2) Introducción. (3) Desarrollo. (4) Conclusiones. (5) Maquetación. (6) Bitácora IA. (7) Carta autor. (8) Revisión. (9) PDF exportado.}

\begin{drawbox}{milcMagenta}{Checklist antes de cerrar}
\checkbox Tema escolar real con destinatario nombrado. \\[1mm] \checkbox 15-20+ respuestas de encuesta reales. \\[1mm] \checkbox Análisis auditado de la IA. \\[1mm] \checkbox Estructura informe técnico completa. \\[1mm] \checkbox Maquetación Docs con 5 elementos profesionales. \\[1mm] \checkbox Bitácora detallada de uso de IA. \\[1mm] \checkbox Carta del autor firmada. \\[1mm] \checkbox PDF exportado y revisado.
\end{drawbox}

\begin{softbox}{milcMagenta}{milcGris}{Plantilla de trabajo}
\textbf{Tema y destinatario.} \textit{Real, nombrado.} \\[1mm] \textbf{Datos.} \textit{15-20+ respuestas con análisis auditado.} \\[1mm] \textbf{Estructura.} \textit{Intro + desarrollo + conclusiones.} \\[1mm] \textbf{Maquetación.} \textit{Docs profesional.} \\[1mm] \textbf{Bitácora IA.} \textit{Modelos + prompts + porcentajes.} \\[1mm] \textbf{Carta.} \textit{Qué + por qué + proceso + compromiso.}
\end{softbox}

\begin{infoband}{milcMagenta}


\textbf{Lo que va al cuaderno (Actividad 3 · CREA + EVALÚA).} Título: «Actividad 3 --- Mini-proyecto integrador». \textbf{Formato:} (a) link al documento; (b) bitácora IA; (c) carta firmada; (d) checklist verificado. \textbf{Sabes que terminaste cuando} podrías enviar el mini-proyecto al destinatario real.
\end{infoband}

\puente{Lo que sigue resume \emph{qué} entregas hoy: mini-proyecto completo + bitácora + carta firmada. El criterio principal: que el mini-proyecto pueda entregarse al destinatario real y tener efecto.}

\milcestacion{milcMostaza}{Producto de la sesión}
\textbf{Mini-proyecto integrador + bitácora IA + carta del autor firmada}.
\begin{softbox}{milcMostaza}{milcGris}{Criterios de calidad}
(1) 6 componentes obligatorios presentes. (2) Datos reales con análisis auditado. (3) Estructura informe técnico completa. (4) Maquetación profesional. (5) Bitácora detallada de IA. (6) Carta del autor firmada.
\end{softbox}

\puente{Antes de cerrar, mira el mini-proyecto desde las cinco dimensiones humanas. Integrar todo el periodo con honestidad es honor del oficio profesional; entregar piezas sueltas es desperdicio de 8 sesiones de trabajo.}

\milcestacion{milcVino}{Evaluación liberadora · cinco dimensiones}

\begin{softbox}{milcVino}{milcGris}{Sentido de la evaluación}
Evaluar no es solo revisar una nota. Es reconocer qué aprendí, qué cambió en mí, qué emociones manejé, cómo me relaciono con la ciudad y la comunidad, y qué le dejaré a quien viene detrás.
\end{softbox}

\textbf{Mi reflexión liberadora en cinco dimensiones.} Completa cada frase iniciada:

\small
\begin{tabularx}{\textwidth}{>{\bfseries\RaggedRight\arraybackslash}p{3.4cm}Y>{\RaggedRight\arraybackslash}p{4.6cm}}
\rowcolor{milcVino}\color{white}Dimensión & \color{white}\bfseries Mi respuesta (frame starter) & \color{white}\bfseries Algo que descubrí\\
1. Personal & Lo que más cambió en mí fue \linead{6.5cm} & \linead{4.2cm}\\
2. Emocional & La emoción más fuerte fue \linead{2.5cm} y la regulé \linead{3cm} & \linead{4.2cm}\\
3. Ciudadana & En democracia esto importa porque \linead{6cm} & \linead{4.2cm}\\
4. Local (Cartago) & En mi barrio sería útil para \linead{6.5cm} & \linead{4.2cm}\\
5. Intergeneracional & A mi abuela/abuelo le diría \linead{6.5cm} & \linead{4.2cm}\\
\end{tabularx}
\normalsize

\begin{infoband}{milcVino}
\textbf{Para la próxima sustentación.} Vuelve a esta tabla en seis meses. Verás que las respuestas cambian — no porque lo aprendido cambie, sino porque tú cambias. La evaluación liberadora es una conversación contigo a través del tiempo.
\end{infoband}


\puente{Y para terminar, tres voces sobre la pieza integradora. Dussel sobre el informe que sostiene voces reales, Marco Aurelio sobre la integración con disciplina, Floridi sobre el mini-proyecto como ética profesional contemporánea.}

\milcestacion{milcGradeDark}{Triángulo de pensamiento · cierre filosófico}

\begin{softbox}{milcVino}{milcGris}{Dussel · lente del nosotros}
\textit{Un mini-proyecto que sostiene voces reales con datos y propuesta es acto profesional político.}

Tu mini-proyecto puede quedarse en archivo escolar o convertirse en propuesta real al destinatario. La filosofía de la liberación pide lo segundo: \textbf{usar la pieza maestra para que algo cambie en tu entorno}. Aunque sea pequeño, un informe entregado a la coordinación o al consejo estudiantil puede generar conversación, decisión, acción. La phronesis del oficio consiste en \textbf{convertir el ejercicio escolar en acto profesional real}.

\textbf{¿Mi mini-proyecto se queda en archivo, o lo entrego al destinatario real para que tenga efecto?}
\end{softbox}
\linead{\linewidth}\\[1mm]
\linead{\linewidth}

\begin{softbox}{milcMostaza}{milcGris}{Estoicismo · Marco Aurelio · cuidado interior}
\textit{Integrar con disciplina es virtud profesional; yuxtaponer sin coherencia es vanidad.}

Es tentador pegar las 8 piezas sin integrar. La sabiduría estoica recomienda lo opuesto: \emph{teje las piezas en un solo argumento}. Esa integración cuesta esfuerzo pero produce piezas profesionales. La phronesis del oficio honra la coherencia sobre la agregación.

\textbf{¿Mi mini-proyecto está tejido como un solo argumento, o solo pegué las 8 piezas en orden?}
\end{softbox}
\linead{\linewidth}\\[1mm]
\linead{\linewidth}

\begin{softbox}{milcTurquesa}{milcGris}{Floridi · lente de la infoesfera}
\textit{El informe profesional con declaración honesta de IA es modelo del oficio del siglo XXI.}

En la era de IA generativa, los informes con declaración honesta y bitácora detallada se posicionan como modelo ético emergente. Tu mini-proyecto de hoy te entrena en esa práctica profesional emergente: \textbf{informes que declaran honestamente uso de IA con porcentajes y prompts}.

\textbf{¿Mi mini-proyecto contribuye al modelo profesional transparente del siglo XXI?}
\end{softbox}
\linead{\linewidth}\\[1mm]
\linead{\linewidth}

\begin{infoband}{milcNegro}
\textbf{Nota para el docente.} Las tres formulaciones del triángulo son la síntesis del autor de la guía, no citas textuales de los pensadores. Si vas a citarlos en clase, remite a la obra original.
\end{infoband}

\milcestacion{milcVino}{Mi autoevaluación · marca una casilla por fila}

{\small
\setlength{\extrarowheight}{2.2pt}
\begin{tabularx}{\textwidth}{>{\RaggedRight\arraybackslash\bfseries}p{3.0cm}YYY}
\rowcolor{milcVino}\color{white}\bfseries Criterio & \color{white}\bfseries Lo logré & \color{white}\bfseries En proceso & \color{white}\bfseries Necesito apoyo\\[.6mm]
Comprensión del tema & \milcopcion{Explico las ideas con mis palabras.} & \milcopcion{Reconozco las ideas, falta precisión.} & \milcopcion{Necesito repasar lo básico.}\\[.8mm]
\rowcolor{milcGris}Pensamiento computacional & \milcopcion{Usé los pilares al armar mi producto.} & \milcopcion{Usé algunos pilares.} & \milcopcion{Necesito releer la estación 2.}\\[.8mm]
Aplicación y producto & \milcopcion{Mi producto cumple los criterios.} & \milcopcion{Está armado, con ajustes pendientes.} & \milcopcion{Aún está en construcción.}\\[.8mm]
\rowcolor{milcGris}Reflexión 5 dimensiones & \milcopcion{Respondí cada una con honestidad.} & \milcopcion{Respondí algunas con detalle.} & \milcopcion{Mis respuestas son muy generales.}\\[.8mm]
Triángulo de pensamiento & \milcopcion{Conecté las 3 voces con mi trabajo.} & \milcopcion{Conecté al menos una.} & \milcopcion{No las relaciono con mi proceso.}\\[.8mm]
\end{tabularx}}

\vspace{3mm}
\textbf{Hoy entendí que:}\\[1mm]
\linead{\linewidth}\\[1mm]
\linead{\linewidth}

\textbf{Mi compromiso para la próxima sesión es:}\\[1mm]
\linead{\linewidth}\\[1mm]
\linead{\linewidth}

\begin{softbox}{milcMostaza}{milcGris}{Cita de despedida · Marco Aurelio}
\textit{``El obstáculo en el camino se vuelve el camino. Lo que estorba al camino se vuelve camino.''}\\[1mm]
Cada error de hoy, bien observado, se convierte en pieza del camino que pisarás mañana.
\end{softbox}

\small
\begin{tabularx}{\textwidth}{>{\bfseries}p{3.6cm}Y}
\rowcolor{milcGradeDark!12}Nombre completo: & \linead{10cm}\\
Fecha: & \linead{4cm} \quad Curso/grupo: \linead{4cm}\\
Mi firma: & \linead{9cm}\\
\end{tabularx}
\normalsize

\begin{infoband}{milcGradeDark}
\textbf{Para la próxima sesión.} Llega con tu mini-proyecto completo, bitácora y carta firmada. En la sesión 10 vas a hacer la \textbf{sustentación pública} del mini-proyecto: 5 minutos defendiendo la pieza ante el grupo. El mini-proyecto de hoy es la pieza; la sustentación de mañana es la defensa pública.
\end{infoband}

\end{document}
